% Part: proof-theory
% Chapter: normalization
% Section: reductions

\documentclass[../../../include/open-logic-section]{subfiles}

\begin{document}

\olfileid[id]{pt}{nor}{red}

\olsection{Konversi Reduksi}

Jika suatu potongan dalam !!a{proof} di \Log{N1i} mempunyai panjang~$1$,
potongan itu merupakan segmen yang terdiri atas satu kemunculan !!{formula},
dan !!{formula} tersebut sekaligus merupakan kesimpulan !!a{introduction}
dan premis utama !!a{elimination}. Dalam bukti normalisasi, potongan seperti
itu ``direduksi'' dengan mengganti sub-!!{proof} yang berakhir dengan
pasangan !!{introduction}/!!{elimination} tersebut dengan !!{proof} baru
yang menghapus pasangan inferensi itu. Penggantian tersebut dapat
menghasilkan potongan baru. Namun, kita dapat meyakinkan diri bahwa
potongan-potongan baru mempunyai peringkat lebih rendah daripada potongan
yang direduksi. Jadi, !!{proof} baru mempunyai panjang potongan lebih rendah
(karena potongan maksimal berpanjang~$1$ telah dihapus) atau peringkat
potongan !!{proof} itu lebih rendah (jika potongan maksimal tersebut merupakan
satu-satunya potongan maksimal).

Sebagai contoh, jika inferensi yang terlibat ialah \Intro{\lif} dan
\Elim{\lif}, $!A \ident !B \lif !C$, dan terdapat
sub-!!{proof}~$\delta_1$ yang berakhir dengan:
\[
\AxiomC{$\Discharge{!B}{x}$}
\RightLabel{$\delta_2$}
\DeduceC{$!C$}
\DischargeRule{\Intro{\lif}}{x}
\UnaryInfC{$!B \lif !C$}
\AxiomC{}
\RightLabel{$\delta_3$}
\DeduceC{$!B$}
\RightLabel{\Elim{\lif}}
\BinaryInfC{$!C$}
\DisplayProof
\]
Untuk menghapus potongan ini, kita mengganti subbukti~$\delta_1$ dalam
$\delta$ dengan yang berikut:
\[
\AxiomC{}
\RightLabel{$\delta_3$}
\DeduceC{$!B$}
\RightLabel{$\Subst{\delta_2}{\delta_3}{!B^x}$}
\DeduceC{$!C$}
\DisplayProof
\]
Dengan $\Subst{\delta_2}{\delta_3}{!B^x}$ kita maksudkan !!{proof} yang
diperoleh dari~$\delta_2$ dengan mengganti setiap asumsi~$!B$ berlabel~$x$
dengan seluruh !!{proof}~$\delta_3$. !!^a{proof} ini tidak lagi mempunyai
asumsi terbuka berbentuk~$!B$ dengan label~$x$. Tidak ada asumsi lain yang
berlabel~$x$, sebab seperti biasa kita mengandaikan bahwa tidak ada dua
!!{formula}s yang muncul sebagai asumsi dengan label sama. Oleh karena itu,
!!{proof} baru mempunyai asumsi terbuka yang sama dengan~$\delta_1$; setiap
inferensi dalam~$\delta_2$ yang melepaskan asumsi melepaskan asumsi yang
sama dalam~$\Subst{\delta_2}{\delta_3}{!B^x}$; dan setiap inferensi dalam
$\delta_3$ yang melepaskan asumsi dapat menghasilkan beberapa salinan yang
melepaskan asumsi bersesuaian dalam~$\Subst{\delta_2}{\delta_3}{!B^x}$.

Setiap potongan yang seluruhnya terletak dalam $\delta_2$ atau $\delta_3$,
maupun seluruhnya di luar~$\delta_1$ dalam~$\delta$, tidak berubah: potongan
itu masih melibatkan !!{formula}s yang sama dan melewati inferensi yang sama,
sehingga khususnya mempunyai peringkat dan panjang sama dengan potongan
bersesuaian dalam~$\delta$. Kita telah menghapus suatu potongan maksimal
berpanjang~$1$ dan karena itu telah mengurangi panjang potongan, atau
mengurangi peringkat potongan (jika panjang potongan~$\delta$ adalah~$1$
dan ini merupakan satu-satunya potongan maksimal).

Namun, dua hal dapat terjadi:
\begin{enumerate}
  \item Dengan mengganti semua asumsi $!B^x$ dengan $\delta_3$, kita telah
  menggandakan semua potongan dalam~$\delta_3$. Jika beberapa di antaranya
  merupakan potongan maksimal, panjang potongan !!{proof} baru dapat lebih
  besar daripada panjang potongan !!{proof} lama. Kita harus menjamin agar
  hal itu tidak terjadi. Caranya ialah menerapkan konversi reduksi hanya pada
  potongan maksimal yang \emph{paling kanan}, yaitu potongan sedemikian
  sehingga tidak ada potongan maksimal pada cabang melalui~$\delta$ yang
  berada di sebelah kanan potongan yang direduksi. Jika terdapat potongan
  dalam~$\delta_3$, potongan itu berada di sebelah kanan potongan yang sedang
  direduksi, sehingga kita tidak sedang menangani potongan paling kanan.
  Dengan selalu memilih potongan paling kanan, kita menjamin bahwa kita
  mengurangi panjang potongan !!{proof}, atau bahkan mengurangi peringkat
  potongannya.
  \item Jika suatu asumsi $!B^x$ dalam $\delta_2$ merupakan premis utama
  !!a{elimination}, dan inferensi terakhir~$\delta_3$ adalah \Elim{\lor},
  \Elim{\lexists}, atau !!a{introduction}, kita telah menghasilkan potongan
  baru dengan mencangkokkan~$\delta_3$ pada $\delta_2$ di asumsi tersebut.
  Hal yang tadinya sekadar asumsi dalam~$\delta_2$ kini merupakan potongan
  berpanjang~$1$ (kesimpulan !!a{introduction} dan premis utama
  !!a{elimination}) atau kemunculan !!{formula} terakhir suatu segmen
  potongan (premis utama !!a{elimination} dan kesimpulan \Elim{\lor} atau
  \Elim{\lexists}). Jika $B^x$ merupakan premis minor \Elim{\lor} atau
  \Elim{\lexists}, kemunculan itu sebelumnya sudah merupakan kemunculan
  !!{formula} pertama suatu segmen, dan mungkin suatu segmen potongan. Dalam
  !!{proof} baru, segmen tersebut kini mungkin lebih panjang. Namun,
  perhatikan bahwa !!{formula} yang membentuk potongan baru atau segmen
  potongan yang sudah ada ini ialah~$!B$, dan
  $\depth{!B} < \depth{!B \lif !C}$. Jadi, walaupun kita mungkin telah
  menambahkan potongan atau memperpanjang potongan, tidak satu pun darinya
  merupakan potongan maksimal. Dengan demikian, kita telah mengurangi
  peringkat potongan, atau, jika masih ada potongan berperingkat sama, telah
  mengurangi panjang potongan.
\end{enumerate}

Proses mereduksi potongan berpanjang~$1$ disebut \emph{konversi reduksi}
(atau \emph{konversi jalan memutar}). Pola bagi beberapa pasangan aturan
diberikan dalam \olref{tab:reduction-conversions}. (Pola lainnya dibiarkan
sebagai latihan.)

\begin{table}
\begin{multline*}
\bottomAlignProof
\AxiomC{$\Discharge{!B}{x}$}
\RightLabel{$\delta_2$}
\shortDeduce
\DeduceC{$!C$}
\DischargeRule{\Intro{\lif}}{x}
\UnaryInfC{$!B \lif !C$}
\AxiomC{}
\RightLabel{$\delta_3$}
\shortDeduce
\DeduceC{$!B$}
\RightLabel{\Elim{\lif}}
\BinaryInfC{$!C$}
\DisplayProof
\quad \longrightarrow\quad
\shoveright{\bottomAlignProof
\AxiomC{}
\RightLabel{$\delta_3$}
\shortDeduce
\DeduceC{$!B$}
\RightLabel{$\Subst{\delta_2}{\delta_3}{!B^x}$}
\shortDeduce
\DeduceC{$!C$}
\DisplayProof}
\\[4ex]
\shoveleft{\bottomAlignProof
\AxiomC{}
\RightLabel{$\delta_2$}
\shortDeduce
\DeduceC{$!B$}
\AxiomC{}
\RightLabel{$\delta_3$}
\shortDeduce
\DeduceC{$!C$}
\RightLabel{\Intro{\land}}
\BinaryInfC{$!B \land !C$}
\RightLabel{\Elim{\land}[1]}
\UnaryInfC{$!B$}
\DisplayProof
\quad \longrightarrow\quad
\bottomAlignProof
\AxiomC{}
\RightLabel{$\delta_2$}
\shortDeduce
\DeduceC{$!B$}
\DisplayProof}\qquad
\shoveright{\bottomAlignProof
\AxiomC{}
\RightLabel{$\delta_2$}
\shortDeduce
\DeduceC{$!B$}
\AxiomC{}
\RightLabel{$\delta_3$}
\shortDeduce
\DeduceC{$!C$}
\RightLabel{\Intro{\land}}
\BinaryInfC{$!B \land !C$}
\RightLabel{\Elim{\land}[2]}
\UnaryInfC{$!C$}
\DisplayProof
\quad \longrightarrow\quad
\bottomAlignProof
\AxiomC{}
\RightLabel{$\delta_3$}
\shortDeduce
\DeduceC{$!C$}
\DisplayProof}
\\[4ex]
\bottomAlignProof
\AxiomC{}
\RightLabel{$\delta_2$}
\shortDeduce
\DeduceC{$!B$}
\RightLabel{\Intro{\lor}[1]}
\UnaryInfC{$!B \lor !C$}
\AxiomC{$\Discharge{!B}{x}$}
\RightLabel{$\delta_3$}
\shortDeduce
\DeduceC{$!D$}
\AxiomC{$\Discharge{!C}{y}$}
\RightLabel{$\delta_4$}
\shortDeduce
\DeduceC{$!D$}
\DischargeRule{\Elim{\lor}}{x\,y}
\TrinaryInfC{$!D$}
\DisplayProof
\quad \longrightarrow\quad
\shoveright{\bottomAlignProof
\AxiomC{}
\RightLabel{$\delta_2$}
\shortDeduce
\DeduceC{$!B$}
\RightLabel{$\Subst{\delta_3}{\delta_2}{!B^x}$}
\shortDeduce
\DeduceC{$!D$}
\DisplayProof}
\\[2ex]
\shoveleft{\bottomAlignProof
\AxiomC{}
\RightLabel{$\delta_2$}
\shortDeduce
\DeduceC{$!B(c)$}
\RightLabel{\Intro{\lforall}}
\UnaryInfC{$\lforall[x][!B(x)]$}
\RightLabel{\Elim{\lforall}}
\UnaryInfC{$!B(t)$}
\DisplayProof
\quad \longrightarrow\quad
\bottomAlignProof
\AxiomC{}
\RightLabel{$\Subst{\delta_2}{t}{c}$}
\shortDeduce
\DeduceC{$!B(t)$}
\DisplayProof}
\\[4ex]
\shoveright{\bottomAlignProof
\AxiomC{}
\RightLabel{$\delta_2$}
\shortDeduce
\DeduceC{$!B(t)$}
\RightLabel{\Intro{\lexists}}
\UnaryInfC{$\lexists[x][!B(x)]$}
\AxiomC{$\Discharge{!B(c)}{x}$}
\RightLabel{$\delta_3$}
\shortDeduce
\DeduceC{$!D$}
\DischargeRule{\Elim{\lexists}}{x}
\BinaryInfC{$!D$}
\DisplayProof
\quad \longrightarrow\quad
\bottomAlignProof
\AxiomC{}
\RightLabel{$\delta_2$}
\shortDeduce
\DeduceC{$!B(t)$}
\RightLabel{$\Subst{\Subst{\delta_3}{t}{c}}{\delta_2}{!B(t)^x}$}
\shortDeduce
\DeduceC{$!D$}
\DisplayProof}
\end{multline*}
\caption{Beberapa konversi reduksi bagi \Log{N1i}.}
\ollabel{tab:reduction-conversions}
\end{table}

\begin{prob}
Berikan konversi reduksi bagi \Intro{\lnot} yang diikuti \Elim{\lnot}.
\end{prob}

Masih ada satu kasus yang harus dipertimbangkan. Definisi potongan mencakup
kasus ketika panjang potongan adalah~$1$ dan $!A$ merupakan kesimpulan
\FalseInt{} sekaligus premis utama !!a{elimination}. Kasus ini ditangani
dengan cara yang serupa dengan kasus ketika $!A$ merupakan kesimpulan
!!a{introduction}. Sebagai contoh, ganti
\[
\bottomAlignProof
\AxiomC{}
\RightLabel{$\delta_2$}
\DeduceC{$\lfalse$}
\RightLabel{\FalseInt}
\UnaryInfC{$!B \lif !C$}
\AxiomC{}
\RightLabel{$\delta_3$}
\DeduceC{$!B$}
\RightLabel{\Elim{\lif}}
\BinaryInfC{$!C$}
\DisplayProof
\quad\text{dengan}\quad
\bottomAlignProof
\AxiomC{}
\RightLabel{$\delta_2$}
\DeduceC{$\lfalse$}
\RightLabel{\FalseInt}
\UnaryInfC{$!C$}
\DisplayProof
\]
Kita harus sedikit berhati-hati dalam kasus $!A \ident (!B \lor !C)$ atau
$!A \ident \lexists[x][!B(x)]$. Misalnya, jika kita mengganti
\[
\bottomAlignProof
\AxiomC{}
\RightLabel{$\delta_2$}
\DeduceC{$\lfalse$}
\RightLabel{\FalseInt}
\UnaryInfC{$!B \lor !C$}
\AxiomC{$\Discharge{!B}{x}$}
\RightLabel{$\delta_3$}
\DeduceC{$!D$}
\AxiomC{$\Discharge{!C}{y}$}
\RightLabel{$\delta_4$}
\DeduceC{$!D$}
\DischargeRule{\Elim{\lor}}{x\,y}
\TrinaryInfC{$!D$}
\DisplayProof
\quad\text{dengan}\quad
\bottomAlignProof
\AxiomC{}
\RightLabel{$\delta_2$}
\DeduceC{$\lfalse$}
\RightLabel{\FalseInt}
\UnaryInfC{$!D$}
\DisplayProof
\]
kita mungkin memperkenalkan potongan baru dengan !!{formula} potongan~$!D$,
dan tidak ada jaminan bahwa
$\depth{!D} < \depth{!B \lor !C}$!{} Jadi, kita harus menempuh cara yang
sama seperti pada reduksi \Intro{\lor}/\Elim{\lor}, yakni menggantinya
dengan
\[
\bottomAlignProof
\AxiomC{}
\RightLabel{$\delta_2$}
\DeduceC{$\lfalse$}
\RightLabel{\FalseInt}
\UnaryInfC{$!B$}
\RightLabel{$\delta_3$}
\DeduceC{$!D$}
\DisplayProof
\quad\text{atau}\quad
\bottomAlignProof
\AxiomC{}
\RightLabel{$\delta_2$}
\DeduceC{$\lfalse$}
\RightLabel{\FalseInt}
\UnaryInfC{$!C$}
\RightLabel{$\delta_4$}
\DeduceC{$!D$}
\DisplayProof
\]

\end{document}
